\directlua{t0 = os.clock()}
\documentclass{article}
%\usepackage{fontspec}
%\setmainfont{Libertinus Serif}
%\usepackage{mathtools,unicode-math}
%\setmainfont{Libertinus Math}
\usepackage[a4paper,margin=2cm]{geometry}
\usepackage{tikz}
\usepackage[dvipsnames]{xcolor}
\usepackage{multicol}
\usepackage[normalem]{ulem}
\usepackage{luacode}
\usepackage{tcolorbox}
\tcbuselibrary{listings}
\tcbuselibrary{skins,breakable}
\usepackage{fancyvrb}
\usepackage[colorlinks]{hyperref}


\usepackage{../luahyperbolic}

\lstdefinelanguage{Lua}{
  morekeywords={and,break,do,else,elseif,end,false,for,function,if,in,
    local,nil,not,or,repeat,return,then,true,until,while},
  sensitive=true,
  morecomment=[l]{--},
  morestring=[b]",
  morestring=[b]'
}

\tcbset{
  myluacolors/.style={
    listing options={
      language=Lua,
      tabsize=2,
      basicstyle=\ttfamily\small,
      keywordstyle=\color{blue!60!black}\bfseries,
      commentstyle=\color{gray}\itshape,
      stringstyle=\color{red!50!black},
      emph={hyper},
      emphstyle={\color{green!40!black}\bfseries},
      showstringspaces=false,
      breaklines=true,
      postbreak=\mbox{\textcolor{red}{$\hookrightarrow$}\space}
    }
  }
}


\newtcblisting{codeandresult}[1]{%
myluacolors,
colback=white,
colframe=black,
fonttitle=\bfseries,
listing side text,
lefthand width=0.62\textwidth,
righthand width=0.35\textwidth,
title=#1
}


\newtcbinputlisting{codefromfile}[2][]{%
  myluacolors,           % shared coloring options
  colback=white,
  colframe=black,
  fonttitle=\bfseries,
  listing file=#2,       % path to external file
  title=#1,
  breakable
}

\title{Package luahyperbolic}
\author{Damien Mégy}
\begin{document}
\maketitle
\begin{center}
\begin{luacode*}
t1 = os.clock()
texio.write_nl(string.format("[TIMING]  %.2f s: beginning intro picture", os.clock() - t0))
local A, B, C = hyper.fundamentalRightTriangle(4,5)
local phi = hyper.automorphism(C,0)
A,B,C = phi(A), phi(B), phi(C)
local geodesics = {
	{hyper.endpoints(A,B)},
	{hyper.endpoints(B,C)},
	{hyper.endpoints(C,A)}
}
geodesics = hyper.propagate_geodesics(geodesics, 24)
local shapes={}
globalNbGeodesics = 0
for k,pair in ipairs(geodesics) do
	globalNbGeodesics = globalNbGeodesics + 1
	table.insert(shapes,hyper.shape_closed_line(pair[1], pair[2]))
end
shapes_string = table.concat(shapes, " ")

texio.write_nl(string.format("[TIMING] %.2f s: computed %d geodesics", os.clock() - t0, globalNbGeodesics))
hyper.tikzBegin("scale=6")
hyper.tikzPrintf(
	"\\fill[even odd rule] (0,0) circle (1) %s ;",
	shapes_string)
hyper.drawRegularPolygon(complex.ZERO, B,5,  "line width=.15cm, cyan")
hyper.tikzEnd()
\end{luacode*}
\end{center}

\begin{luacode*}
globalTimeIntroPicture = os.clock() - t1
texio.write_nl(string.format("[TIMING] %.2f s: finished tikz picture in %.2f  seconds", os.clock() - t0, globalTimeIntroPicture))
\end{luacode*}

\begin{multicols}{2}
\tableofcontents
\end{multicols}



\newpage
\section{Introduction}

\subsection{What you can do with this package}
\begin{enumerate}
\item Draw hyperbolic geometry pictures, high precision (pdf output).
\item Use an easy syntax. At least, easier than with \LaTeX macros and blody pgf.
\item Export generated pictures to external tikz files if you want.
\item Read the source code and learn hyperbolic geometry.
\item Compile **without** installing python packages or lua packages etc.
\item Compile **without**  "shell-escape" mode.
\item Compile everything in a single pass. No (Leeloo Dallas) multipass. Want to illustrate Euclid's fifth postulate ? Do it right there in your tex file without having to include external files :

\begin{codeandresult}{Euclid's 5th Postulate}
\begin{luacode*}
hyper.tikzBegin("scale=2.5")
local P = complex(0.5,-0.2)
local A = complex.exp_i(math.pi/10)
for k=1,5 do
	hyper.drawLine(P, A^k, "teal")
end
hyper.labelPoint(P, "$P$", "left=.2cm")
local style = "very thick, dashed, red"
hyper.drawLine(complex.J,-complex.I,style)
hyper.tikzEnd("euclid5thPostulate.tikz")
\end{luacode*}
\end{codeandresult}

\end{enumerate}

\subsection{What this package cannot do for you}

\begin{enumerate}
\item Really intensive computations, millions of geodesics etc.  Source code is not optimized to keep maximum readability. All computations involve automorphisms. See other remarks below.
\item Coffee
\end{enumerate}

\subsection{Efficiency}

For heavy computations, please use a \sout{Python} C library. Nevertheless, this library is still reasonably efficient and usable:
\begin{itemize}
\item The first tiling illustration of this documentation (hyperbolic 2-4-5 tiling with \directlua{tex.print(globalNbGeodesics)} geodesics) compiled in \directlua{tex.print(math.round(100*globalTimeIntroPicture)/100)} seconds on a 2020 intel macbook air laptop.
\item This entire documentation including all pictures compiles in approx. 15 seconds on the same machine, the majority of the time weing spent on the tilings. The package is suitable for producing lecture notes in hyperbolic geometry, which was the reason for all of this. For an entire book with dozens of pictures, you may need to externalize some of the pictures.
\item Tests show that the bottleneck is TikZ, not Lua. In a future version, some future maintainer may add other fontends to the library, including drawing directly to the pdf with pdf primitives wich are several orders of magnitude faster than tikz.
\end{itemize}


\subsection{How to load the package}

\subsubsection*{Local install}
Simply put the \verb+luahyperbolic.sty+ file in your working directory and  add the following line to your tex preamble :
\begin{verbatim}
\usepackage{luahyperbolic}
\end{verbatim}

The file \texttt{luahyperbolic.sty} \textbf{contains everything needed}, including the necessary library for complex numbers.  It can be sent to colleagues by email etc.

It will declare two lua globals : \texttt{complex} and \texttt{hyper}, and define a \texttt{hyperbolicdisk} environment for convenience.

If the sty file is not in the same directory as your tex file, you can write \verb+\usepackage{../luahyperbolic}+


\subsubsection*{CTAN} Not  available yet.

\subsection{How to customize the package}

The \texttt{luahyperbolic.sty}  is generated from three lua modules : \texttt{complex.lua}, \texttt{luahyperbolic-core.lua} and \texttt{luahyperbolic-tikz.lua}. These files are in the \texttt{lua/} subdirectory of the package repo. 

To build the \texttt{luahyperbolic.sty} from lua source files, execute \texttt{lua build.lua} in the root directory of the package. This will backup the old sty file to the \texttt{backup/} directory, then buid and replace the existing sty file.

\subsection{Environment hyperbolicdisk}
The package provides the latex environment \verb+\begin{hyperbolicdisk} ... \end{hyperbolicdisk}+ which is defined as follows:

\begin{verbatim}
\newenvironment{hyperbolicdisk}[1][scale=3]{%
  \begin{tikzpicture}[#1]
    \draw[very thick] (0,0) circle (1);
    \clip (0,0) circle (1);
}{%
  \end{tikzpicture}
}
\end{verbatim}

In this documentation, we prefer to use the lua functions 
 \verb+hyper.tikzBegin(options)+ and \verb+hyper.tikzEnd(export_file)+, that enable export functionalities, see below.
 
\subsubsection{Exporting the pictures}

Simply add a filename to the closing \texttt{hyper.tikzEnd()} instruction, as below:
\begin{center}
\texttt{hyper.tikzEnd("myfile.tikz")}
\end{center}

This will save the current picture (including \verb+\begin{tikzpicture}[options]+ and \verb+\end{tikzpicture}+


It is also possible to export the current picture at any point with \verb+hyper.tikzExport(optional-filename)+.
If no filename is provided, the package will name it \verb+hyper_picture_N+ with \texttt{N} the picture counter. Luahyperbolic keeps count of how many pictures have been produced in the document.
This allows exporting images in loops, or drawing several steps of the same picture.




%The code and result are displayed side by side thanks to the \verb+tcblisting+ environment of the \verb+tcolorbox+ package


\section{Drawing elements}

% ------------------ Figure 1 : Points ------------------

\subsection{Points}


\begin{codeandresult}{Draw and label points}
\begin{luacode*}
local A = complex(0.2,0.1)
local B = complex(-0.3,0.5)
local C = complex(0.6,-0.5)
local D = complex(-0.6,0.2)
hyper.tikzBegin()
hyper.drawPoint(A)
hyper.drawPoints(B, C, "teal")
hyper.drawPoint(D, "white, draw=red")
hyper.labelPoint(A,"$A$") -- default: above left
hyper.labelPoints(B,"$B$", C, "$C$", "above right, teal")
hyper.labelPoint(D,"$D$","below=.4cm, red")
hyper.tikzEnd()
\end{luacode*}
\end{codeandresult}


\begin{codeandresult}{Random points with uniform density}
\begin{luacode*}
local rMin, rMax = .3, .8
hyper.tikzBegin()
for k=1,100 do
	local M = hyper.randomPoint(rMin,rMax)
	hyper.drawPoint(M)
end
hyper.tikzPrintf("\\draw[dashed] (0,0) circle (%s) ;", rMin)
hyper.tikzPrintf("\\draw[dashed] (0,0) circle (%s) ;", rMax)
hyper.tikzEnd()
\end{luacode*}
\end{codeandresult}

\subsection{Geodesic segments, lines and rays}


\begin{codeandresult}{Geodesic segments}
\begin{luacode*}
hyper.tikzBegin()
local A = complex(0.2,0.1)
local B = complex(-0.3,0.5)
local C = complex(0.6,-0.5)
hyper.drawSegment(A, B, "thick,blue")
hyper.drawSegment(A, C, "dashed,ForestGreen")
hyper.drawPoints(A,B,C)
hyper.tikzEnd()
\end{luacode*}
\end{codeandresult}


\begin{codeandresult}{Lines and rays (through two points)}
\begin{luacode*}
hyper.tikzBegin()
local A = complex(-0.2,-0.5)
local B = complex(0.5,0.6)
local C = complex(-0.5,0.5)
local D = complex(-0.7,-0.3)
hyper.drawLine(A, B, "red")
hyper.drawRay(C,D, "blue")
hyper.drawLine(D,-D, "green!50!black")
hyper.drawLine(A,complex(1,0), "gray")
hyper.drawLine(complex(1,0), complex(0,1), "teal")
hyper.drawPoints(A, B, C, D)
hyper.labelPoints(A, "$A$", B, "$B$", C, "$C$", D, "$D$")
hyper.tikzEnd()
\end{luacode*}
\end{codeandresult}

\begin{codeandresult}{Triangles and polygons}
\begin{luacode*}
hyper.tikzBegin()
hyper.drawTriangle(
	complex(.9,.2),
	complex(.5,.7),
	complex(.7,-.2), "red")
local points={}
for k=1,5 do
	table.insert(points, complex.polar(.5,k*2*math.pi/5))
end
hyper.drawPolygonFromTable(points,"teal")
hyper.tikzEnd()
\end{luacode*}
\end{codeandresult}

\begin{codeandresult}{Regular polygons}
\begin{luacode*}
hyper.tikzBegin()
local O1 = complex(.5,.2)
local A1 = complex(.7,.3)
local O2 = complex(-0.1,0)
local A2 = complex(-.4,.3)
local O3 = complex(0,.7)
local A3 = complex(0,.9)
hyper.drawRegularPolygon(O1, A1, 3, "red")
hyper.drawRegularPolygon(O2, A2, 5, "blue")
hyper.drawRegularPolygon(O3, A3, 8, "green!50!black")
hyper.drawPoints(O1,O2,O3)
hyper.tikzEnd()
\end{luacode*}
\end{codeandresult}


\begin{codeandresult}{Polygonal lines}
\begin{luacode*}
hyper.tikzBegin()
local A = complex(-.9,.2)
local B = complex(-.5,.8)
local C = complex(-.7,-.4)
local D = complex(.7,-.5)
hyper.drawPolyline(A, B, C, D, "red")
local points={}
for k=1,6 do
	table.insert(points, hyper.randomPoint(0,.7))
end
hyper.drawPolylineFromTable(points,"teal")
hyper.tikzEnd()
\end{luacode*}
\end{codeandresult}




\begin{codeandresult}{Lines and rays (point and tangent vector)}
\begin{luacode*}
hyper.tikzBegin()
local A = complex(-0.2,-0.5)
local B = complex(-0.5,0.5)
local C = complex(0,0.2)
local v =  complex(1,0)
hyper.drawLineFromVector(A, v, "red")
hyper.drawRayFromVector(B, v, "blue")
hyper.drawRayFromVector(C, v, "Green")
hyper.drawPoints(A, B,C)
hyper.drawTangentVector(A,v,"red")
hyper.drawTangentVector(B,v,"blue")
hyper.drawTangentVector(C,v,"Green")
hyper.tikzEnd()
\end{luacode*}
\end{codeandresult}






\begin{codeandresult}{Perpendicular bisector}
\begin{luacode*}
hyper.tikzBegin()
local A = complex(-0.2,-0.5)
local B = complex(-0.7,0.3)
hyper.drawPerpendicularBisector(A,B,"red")
hyper.drawPoints(A,B)
hyper.drawSegment(A, B, "teal")
hyper.labelPoints(A,"$A$", B, "$B$", "above right")
hyper.tikzEnd()
\end{luacode*}
\end{codeandresult}


\begin{codeandresult}{Angle bisector}
\begin{luacode*}
hyper.tikzBegin()
local A = complex(-0.2,-0.6)
local O = complex(0.3,0.3)
local B = complex(-0.7,0.3)
hyper.drawAngleBisector(A,O,B,"red")
hyper.drawPoints(A,O,B)
hyper.drawPolyline(A,O,B, "teal")
hyper.labelPoints(A, "$A$", O, "$O$", B, "$B$")
hyper.tikzEnd()
\end{luacode*}
\end{codeandresult}

\subsection{Circles, semicircles, arcs}

\begin{codeandresult}{Center and radius}
\begin{luacode*}
hyper.tikzBegin()
local points = {
	complex(0,0),
	complex(-0.5,-0.4),
	complex(0.8,0.1),
	complex(0,0.9)
}
for _,point in ipairs(points) do
	hyper.drawCircle(point,1,"teal")
	hyper.drawPoint(point)
end
-- NOTE : all circles have same radius=1 !
hyper.tikzEnd()
\end{luacode*}
\end{codeandresult}


\begin{codeandresult}{Circle through point}
\begin{luacode*}
hyper.tikzBegin()
local O = complex(0.3,0.2)
local P = complex(0.7,-0.4)
local Q = complex(-0.1,0.4)
hyper.drawCircleThrough(O, P, "red")
-- or by hand : 
hyper.drawCircle(O, hyper.distance(O,Q), "blue")
hyper.drawPoints(O, P, Q)
hyper.labelPoints(O, "$O$", P, "$P$", Q, "$Q$")
hyper.tikzEnd()
\end{luacode*}
\end{codeandresult}

\begin{codeandresult}{Semicircles}
\begin{luacode*}
hyper.tikzBegin()
local O = complex(0,0)
local A = complex(-0.3,0)
local P = complex(0.5,0.2)
local B = complex(0.7,-0.4)
hyper.drawSemicircle(O, A, "blue")
hyper.drawSemicircle(P, B, "red")
hyper.drawPoints(O, A, P, B)
hyper.labelPoints(O, "$O$", A, "$A$", P, "$P$", B, "$B$")
hyper.drawLine(P,B, "gray, dashed")
hyper.tikzEnd()
\end{luacode*}
\end{codeandresult}


\begin{codeandresult}{Arc (center and endpoints)}
\begin{luacode*}
hyper.tikzBegin()
local O = complex(0,0)
local A = complex(0.5,0)
local B = complex.polar(0.5,math.pi/3)
hyper.drawArc(O,A,B, "red")
hyper.drawPoints(O,A,B)
hyper.tikzEnd()
\end{luacode*}
\end{codeandresult}

\begin{codeandresult}{Draw angle $\widehat{AOB}$}
\begin{luacode*}
hyper.tikzBegin()
local O = complex(-.5,-.3)
local A = complex(0.5,0)
local B = complex(0.2,.7)
hyper.drawLines(O, A, O, B, "teal")
hyper.drawAngle(A,O,B, "red, very thick")
hyper.drawPoints(A,O,B)
hyper.labelPoints(A, "$A$", O, "$O$", B, "$B$")
-- poor man's labelAngle:
hyper.labelPoint(O,"$\\theta$", "above=.5cm, right=.6cm")
hyper.tikzEnd()
\end{luacode*}
\end{codeandresult}

\begin{codeandresult}{Tangent at point}
\begin{luacode*}
hyper.tikzBegin()
local O = complex(0.3,0.2)
local P = complex(0.1,-0.4)
local Q = hyper.rotation(O,math.pi/3)(P)
hyper.drawCircleThrough(O, P)
hyper.drawPoints(O, P, Q)
hyper.drawTangentAt(O,P,"red")
hyper.drawTangentAt(O,Q,"teal")
hyper.tikzEnd()
\end{luacode*}
\end{codeandresult}

%%%%%%%%%%%%%%%%%%%%%%%%%%%%%%%%%%%%%%%%%%%%
\subsection{Horocycles}

\begin{codeandresult}{Horocycles}
\begin{luacode*}
hyper.tikzBegin()
local theta=math.pi/3
local idealPoint = complex.exp_i(theta)
local N = 10
for k=1,N-1 do
	local point = idealPoint * (2*k/N-1)
	hyper.drawHorocycle(idealPoint,point,"red")
	-- and an orthogonal geodesic : 
	local endPoint = complex.exp_i(theta+2*k*math.pi/N)
	hyper.drawLine(idealPoint,endPoint,"gray")
end
hyper.tikzEnd()
\end{luacode*}
\end{codeandresult}



\section{Defining points}

\subsection{Midpoint, interpolate, barycenter}

\begin{codeandresult}{Midpoint}
\begin{luacode*}
hyper.tikzBegin()
local A = complex(0.4,0.4)
local B = complex(-0.8,-0.3)
local M = hyper.midpoint(A,B)
hyper.drawSegment(A,B, "teal")
hyper.drawPoints(A,B,M)
hyper.labelPoints(A, "$A$", B, "$B$", M, "$M$")
hyper.tikzEnd()
\end{luacode*}
\end{codeandresult}

\begin{codeandresult}{Interpolate and barycenter}
\begin{luacode*}
hyper.tikzBegin()
local A = complex(0.3,0.8)
local B = complex(-0.7,-0.4)
hyper.drawSegment(A,B, "teal")
for k=0,10 do
	local M = hyper.interpolate(A, B, k/10)
	hyper.drawPoint(M)
	if k % 2 == 0 then
		hyper.labelPoint(M, "$M_{" .. k .. "}$", "below right")
	end
end
hyper.tikzEnd()
\end{luacode*}
\end{codeandresult}

\subsection{Intersections}

\begin{codeandresult}{Line-Line (LL) intersection}
\begin{luacode*}
hyper.tikzBegin()
local A = complex(0.4,0.3)
local B = complex(-0.4,-0.7)
local C = complex(0,0.7)
local D = complex(0.8,-0.2)
local P = hyper.interLL(A, B, C, D)
hyper.drawLines(A, B, C, D, "teal")
hyper.drawPoints(A, B, C, D,"teal")
hyper.drawPoint(P,"red")
hyper.labelPoint(A, "$A$", "above")
hyper.labelPoint(B, "$B$")
hyper.labelPoint(C, "$C$", "left")
hyper.labelPoint(D, "$D$", "above")
hyper.labelPoint(P, "$P$", "left=.1cm, red")
hyper.tikzEnd()
\end{luacode*}
\end{codeandresult}



\begin{codeandresult}{Line-circle (LC) intersection}
\begin{luacode*}
hyper.tikzBegin()
local A = complex(0.5,0.4)
local B = complex(-0.3,-0.7)
local O = complex(-0.2,0.1)
local radius = 1
local P, Q = hyper.interLC(A, B, O, radius)
hyper.drawLine(A, B, "teal")
hyper.drawCircle(O,radius, "teal")
hyper.drawPoints(A, B, O, "teal")
hyper.drawPoints(P, Q, "red")
hyper.labelPoints(A, "$A$", B, "$B$", O, "$O$")
hyper.labelPoints(P, "$P$", Q, "$Q$", "red, below right")
hyper.tikzEnd()
\end{luacode*}
\end{codeandresult}

\begin{codeandresult}{Circle-circle (CC) intersection}
\begin{luacode*}
hyper.tikzBegin()
local O1 = complex(-0.6,-.3)
local O2 = complex(0.5,0)
local r1, r2 = 1, 2
local P, Q = hyper.interCC(O1, r1, O2, r2)
hyper.drawCircle(O1,r1, "teal")
hyper.drawCircle(O2,r2, "teal")
hyper.drawPoints(P, Q, "red")
hyper.drawPoints(O1, O2, "teal")
hyper.labelPoints(O1, "$O_1$", O2, "$O_2$", "left, teal")
hyper.labelPoints(P, "$P$", "above left, red")
hyper.labelPoint(Q, "$Q$", "red, below")
hyper.tikzEnd()
\end{luacode*}
\end{codeandresult}



\subsection{Reflections and projections}


\begin{codeandresult}{Reflection across a geodesic line}
\begin{luacode*}
hyper.tikzBegin()
local A = complex(0.8,0.2)
local B = complex(-0.3,0.5)
local M = complex(0.3,-0.6)
local Mp = (hyper.reflection(A,B))(M)
hyper.drawLine(A, B, "teal")
hyper.drawPoints(A, B, M)
hyper.drawPoint(Mp,"red")
hyper.drawSegment(M,Mp,"dashed,red")
hyper.labelPoint(Mp,"$M'$", "left, red")
hyper.labelPoints(A, "$A$", B, "$B$", M, "$M$", "below")
hyper.tikzEnd()
\end{luacode*}
\end{codeandresult}

\begin{codeandresult}{Projection onto a geodesic line}
\begin{luacode*}
hyper.tikzBegin()
local A = complex(0.8,0.2)
local B = complex(-0.3,0.5)
local M = complex(0.3,-0.6)
local P = (hyper.projection(A,B))(M)
hyper.drawLine(A, B, "teal")
hyper.drawPoints(A, B, M)
hyper.drawPoint(P,"red")
hyper.drawSegment(M,P,"dashed,red")
hyper.labelPoint(P,"$P$", "above, red")
hyper.labelPoints(A, "$A$", B, "$B$", M, "$M$", "below")
hyper.tikzEnd()
\end{luacode*}
\end{codeandresult}

\subsection{(Elliptic) Rotations}


\begin{codeandresult}{(Elliptic) rotation around a point (regular octogon)}
\begin{luacode*}
hyper.tikzBegin()
local O = complex(0.3,0.2)
local M = complex(0.5,-0.4)
hyper.drawPoints(O)
for k=1,8 do
	local P = (hyper.rotation(O,2*math.pi/8))(M)
	hyper.drawSegment(M, P, "red")
	hyper.drawSegment(O,M, "dashed,red")
	M = P
end
hyper.drawCircle(O,hyper.distance(O,M),"thin,gray")
hyper.tikzEnd()
\end{luacode*}
\end{codeandresult}

\begin{codeandresult}{Symmetry around a point (rotation by $180^\circ$)}
\begin{luacode*}
hyper.tikzBegin()
local O = complex(0.2,0.2)
local s = hyper.symmetry(O)
hyper.drawPoints(O)
local points={complex(-.8,.1), complex(.6,0), complex(0,-.5)}
for k, P in ipairs(points) do
	hyper.drawCircleThrough(O, P, "gray!10")
	local sP = s(P)
	hyper.drawPoints(P,sP)
	hyper.drawSegment(P,sP, "red, dashed")
	hyper.labelPoint(P, "$P_{".. k .."}$")
	hyper.labelPoint(sP, "$\\sigma(P_{".. k .."})$")
end
hyper.tikzEnd()
\end{luacode*}
\end{codeandresult}


\begin{codeandresult}{Rotation by $60^\circ$}
\begin{luacode*}
hyper.tikzBegin()
local O = complex(0.1,0.1)
local r = hyper.rotation(O,math.pi/3)
hyper.drawPoints(O)
for k=1,5 do
	local P = hyper.randomPoint(0.3,0.9)
	local rP = r(P)
	hyper.drawCircleThrough(O,P, "gray!10")
	hyper.drawArc(O,P,rP,"red")
	hyper.drawPolyline(P,O,rP, "red!10")
	hyper.drawPoints(P,rP)
end
hyper.tikzEnd()
\end{luacode*}
\end{codeandresult}


\subsection{Orbits under automorphisms}


\begin{codeandresult}{Orbit under elliptic automorphism (rotation)}
\begin{luacode*}
hyper.tikzBegin()
local nbPoints = 7
local nbIterations = 5
local O = complex(0.1,0.2)
local rot = hyper.rotation(O,math.pi/20) -- angle in radians
hyper.drawPoint(O,"red")
for k=1,nbPoints do
	local P = hyper.randomPoint(0,.9)
	hyper.drawCircleThrough(O,P,"draw opacity=0.2")
	hyper.drawPointOrbit(P, rot, nbIterations) 
end
hyper.tikzEnd()
\end{luacode*}
\end{codeandresult}


\begin{codeandresult}{Orbit under hyperbolic automorphism}
\begin{luacode*}
hyper.tikzBegin()
local nbPoints = 10
local nbIterations = 5
local A = complex(0.1,0.1)
local phi = hyper.automorphism(A,0)
for k=1,nbPoints do
	local P = hyper.randomPoint(0,.9)
	hyper.drawPointOrbit(P, phi, nbIterations) 
end
hyper.drawPoint(A,"red")
hyper.labelPoint(A,"$A$", "above left, red")
hyper.drawPointOrbit(A, phi, nbIterations, "red") 
hyper.tikzEnd()
\end{luacode*}
\end{codeandresult}


\begin{codeandresult}{Orbit under parabolic automorphism}
\begin{luacode*}
hyper.tikzBegin()
local nbIterations = 7
limit = complex.exp_i(math.pi/3) -- on circle
local phi = hyper.parabolic(limit,math.pi/20)
local points={
	complex(0,0), complex(0,.7),
	complex(.4,.4), complex(.3,.7),
	complex(-.8,.4), complex(-.4,0),
	complex(.4,-.45)}
for _,P in ipairs(points) do
	hyper.drawPointOrbit(P, phi, nbIterations) 
	hyper.drawHorocycle(limit, P, "red, draw opacity=.3")
end
hyper.tikzEnd()
\end{luacode*}
\end{codeandresult}


\subsection{Triangle geometry (CAUTION)}

\begin{codeandresult}{Centroid (intersection of medians)}
\begin{luacode*}
hyper.tikzBegin()
local A = complex(0.1,0.8)
local B = complex(-0.7,-0.3)
local C = complex(0.6,-0.2)
local AA = hyper.midpoint(B,C)
local BB = hyper.midpoint(C,A)
local CC = hyper.midpoint(A,B)
local G = hyper.triangleCentroid(A,B,C) 
hyper.drawTriangle(A,B,C)
hyper.drawSegments(A, AA, B, BB, C, CC, "teal")
hyper.drawPoints(AA,BB,CC)
hyper.drawPoint(G,"white, draw=teal")
-- CAUTION : 
local wrong = hyper.interpolate(A,AA,2/3)
hyper.drawPoint(wrong, "red")
hyper.tikzEnd()
\end{luacode*}
\end{codeandresult}


\begin{codeandresult}{Incenter and incircle}
\begin{luacode*}
hyper.tikzBegin()
local A = complex(0.1,0.8)
local B = complex(-0.7,-0.3)
local C = complex(0.6,-0.2)
local I = hyper.triangleIncenter(A,B,C)
hyper.drawTriangle(A,B,C,"teal")
hyper.drawPoint(I,"white, draw=red")
hyper.labelPoint(I,"$I$", "below left")
hyper.drawIncircle(A,B,C,"red")
-- show intersection : 
local mystyle = "dashed, red!50"
hyper.drawAngleBisector(A,B, C, mystyle)
hyper.drawAngleBisector(B,C,A, mystyle)
hyper.drawAngleBisector(C,A,B, mystyle)
hyper.tikzEnd()
\end{luacode*}
\end{codeandresult}

\begin{codeandresult}{Circumcenter and circumcircle (CAUTION : does not always exist)}
\begin{luacode*}
hyper.tikzBegin()
local A = complex(0.1,0.8)
local B = complex(-0.6,-0.5)
local C = complex(0.6,-0.2)
hyper.drawTriangle(A,B,C,"teal")
local mystyle = "dashed, red!50"
hyper.drawPerpendicularBisector(A, B, mystyle)
hyper.drawPerpendicularBisector(B, C, mystyle)
hyper.drawPerpendicularBisector(C, A, mystyle)
local O = hyper.triangleCircumcenter(A,B,C)
hyper.drawCircumcircle(A,B,C,"red")
hyper.labelPoint(O,"$O$", "below left")
local BC = hyper.midpoint(B,C)
local CA = hyper.midpoint(C,A)
local AB = hyper.midpoint(A,B)
hyper.drawRightAngle(O, CA, C,"red",1/8)
hyper.drawRightAngle(O,BC,B,"red",1/8)
hyper.drawPoints(O,A,B,C,AB,BC,CA)
hyper.labelPoints(A, "$A$", B, "$B$", C, "$C$","above")
hyper.tikzEnd()
\end{luacode*}
\end{codeandresult}


\begin{codeandresult}{Orthocenter (CAUTION : does not always exist)}
\begin{luacode*}
hyper.tikzBegin()
local A = complex(0.1,0.8)
local B = complex(-0.7,-0.4)
local C = complex(0.6,-0.2)
hyper.drawTriangle(A,B,C,"teal")
local mystyle = "dashed, red!50"
hyper.drawPerpendicularThrough(A,B,C,mystyle)
hyper.drawPerpendicularThrough(B,C,A,mystyle)
hyper.drawPerpendicularThrough(C,A,B,mystyle)
local AA = hyper.projection(B,C)(A)
local BB = hyper.projection(C,A)(B)
local CC = hyper.projection(A,B)(C)
local style2 = "very thick, red"
hyper.drawRightAngle(A,AA,B,style2,1/15)
hyper.drawRightAngle(B,BB,C,style2,1/8)
hyper.drawRightAngle(C,CC,A,style2,1/12)
hyper.labelPoints(A, "$A$", B, "$B$", C, "$C$")
local H = hyper.triangleOrthocenter(A,B,C)
hyper.drawPoints(AA,BB,CC,H,"above")
hyper.tikzEnd()
\end{luacode*}
\end{codeandresult}



\newpage

\section{Tilings}
The two following tilings  are drawn with \verb+TILING_DEPTH=12+
\directlua{TILING_DEPTH = 12}

\begin{codeandresult}{Triangle tiling}
\begin{luacode*}
hyper.tikzBegin()
local A, B, C = hyper.fundamentalRightTriangle(4,5)
local geodesics = {
{hyper.endpoints(A,B)},
	{hyper.endpoints(B,C)},
	{hyper.endpoints(C,A)}}
geodesics = hyper.propagate_geodesics(geodesics, TILING_DEPTH)
hyper.drawLinesFromTable(geodesics, "green!50!black, draw opacity=.5")
hyper.drawTriangle(A,B,C,"red, very thick")
hyper.tikzEnd()
\end{luacode*}
\end{codeandresult}

\begin{codeandresult}{Triangle tiling with one angle zero}
\begin{luacode*}
hyper.tikzBegin()
local A, B, C = hyper.fundamentalIdealTriangle(3,2)
local geodesics = {
	{hyper.endpoints(A,B)},
	{hyper.endpoints(B,C)},
	{hyper.endpoints(C,A)}
}
geodesics = hyper.propagate_geodesics(geodesics, TILING_DEPTH)
hyper.drawLinesFromTable(geodesics, "teal, draw opacity=.5")
hyper.drawTriangle(A,B,C,"red, very thick")
hyper.tikzEnd()
\end{luacode*}
\end{codeandresult}


\subsection{Source code of first picture}

\begin{lstlisting}[
	language=Lua,
      tabsize=2,
      basicstyle=\ttfamily\small,
      keywordstyle=\color{blue!60!black}\bfseries,
      commentstyle=\color{gray}\itshape,
      stringstyle=\color{red!50!black},
      emph={hyper},
      emphstyle={\color{green!40!black}\bfseries},
      showstringspaces=false,
      breaklines=true,
      postbreak=\mbox{\textcolor{red}{$\hookrightarrow$}\space}]

\begin{luacode*}
local A, B, C = hyper.fundamentalRightTriangle(4,5)
local phi = hyper.automorphism(C,0)
A,B,C = phi(A), phi(B), phi(C) -- center on pentagon
local geodesics = {
	{hyper.endpoints(A,B)},
	{hyper.endpoints(B,C)},
	{hyper.endpoints(C,A)}
}
geodesics = hyper.propagate_geodesics(geodesics, 24)
local shapes={}
for k,pair in ipairs(geodesics) do
	table.insert(shapes,hyper.shape_closed_line(pair[1], pair[2]))
end
shapes_string = table.concat(shapes, " ")
hyper.tikzBegin("scale=6")
hyper.tikzPrintf("\\fill[even odd rule] (0,0) circle (1) %s ;",shapes_string)
hyper.drawRegularPolygon(complex.ZERO, B,5,  "line width=.15cm, cyan")
hyper.tikzEnd()
\end{luacode*}
\end{lstlisting}

\subsection{More images just for fun}

\begin{luacode*}
local N=5
local DEPTH=8
for k=2,N do
	for l=k, N do
		if 1/k+1/l <1 then
			local A, B, C = hyper.fundamentalIdealTriangle(k,l)
			local geodesics = {
				{hyper.endpoints(A,B)},
				{hyper.endpoints(B,C)},
				{hyper.endpoints(C,A)}
			}
			geodesics = hyper.propagate_geodesics(geodesics, DEPTH)
			local shapes={}
			for k,pair in ipairs(geodesics) do
				table.insert(shapes,hyper.shape_closed_line(pair[1], pair[2]))
			end
			shapes_string = table.concat(shapes, " ")
			hyper.tikzBegin("scale=2.7")
			hyper.tikzPrintf("\\fill[even odd rule](0,0) circle (1) %s ; ", shapes_string)
			hyper.tikzEnd()
			tex.print("~")
		end
	end
end 
\end{luacode*}

\section{Overview of the lua libraries}

The Lua files are in the \verb+lua/+ subdirectory of the repo.

To build the \verb+luahyperbolic.sty+ from the lua source files, run \verb+lua build.lua+ on the root directory. (Lua must be installed)

The Lua source files are not needed to compile a latex document using luahyperbolic, only the sty file is necessary.

\begin{luacode*}
-- Print all entries of a given type in a module
function print_module_type(module, wanted_type, sep)
    sep = sep or ", "
    local names = {}
    for k, v in pairs(module) do
        if type(v) == wanted_type then
            table.insert(names, k)
        end
    end
    table.sort(names)
    for i, name in ipairs(names) do
        tex.sprint(-2, name)
        tex.sprint(sep)
    end

end
\end{luacode*}

\subsection{complex.lua}
\subsubsection*{Functions}
\begin{multicols}{6}
\noindent\directlua{print_module_type(complex, "function", "\\\\")}
\end{multicols}
\subsubsection*{Constants (numbers and complex numbers)}
\directlua{print_module_type(complex, "number", ", ")}
\directlua{print_module_type(complex, "table", ", ")}



\subsection{luahyperbolic-tikz.lua and hyperbolic-core}
\subsubsection*{Functions}
\begin{multicols}{3}
\noindent\directlua{print_module_type(hyper, "function", "\\\\")}
\end{multicols}
\subsubsection*{Constants}
\textbf{Numbers : }
\directlua{print_module_type(hyper, "number", ", ")}

\textbf{Strings : }
\directlua{print_module_type(hyper, "string", ", ")}

\textbf{Tables : }
\directlua{print_module_type(hyper, "table", ", ")}

\end{document}